\chapter{Mantenimiento}
\label{chap:mantenimiento}

\section{Objetivo del mantenimiento}
\label{sec:mantenimiento-objetivo}

El mantenimiento de la plataforma debe proteger la consistencia del torneo. La prioridad no es cambiar pantallas rapidamente, sino conservar la relacion entre registros, equipos, grupos, batallas, estados, historial, comunicaciones y panel interno. Los modulos con mayor impacto estan descritos en \cref{chap:aplicaciones-django}, \cref{chap:servicios-logica-negocio}, \cref{chap:frontend-templates} y \cref{chap:javascript-ajax}.

Este capitulo sirve como guia para futuros desarrolladores que deban corregir errores, agregar formatos, ajustar comunicaciones o preparar una version para evento.

\section{Modulos sensibles}
\label{sec:mantenimiento-modulos}

\begin{longtable}{p{3.6cm}p{4.8cm}p{4.4cm}}
\caption{Modulos sensibles de mantenimiento.}
\label{tab:mantenimiento-modulos}\\
\toprule
Modulo & Archivos principales & Riesgo de cambio \\
\midrule
Configuracion & \archivo{config/settings.py}, \archivo{.env.example}. & Romper arranque por entorno, seguridad o conexion a base de datos. \\
Registro & \archivo{apps/participants/models.py}, \archivo{forms.py}, \archivo{services.py}, \archivo{views.py}. & Permitir duplicados o no sincronizar equipos confirmados. \\
Torneo & \archivo{apps/tournament/services.py}, \archivo{formats.py}, \archivo{recommendations.py}, \archivo{models.py}. & Inconsistencia en fases, grupos, batallas, estados o historial. \\
Panel interno & \archivo{control_stage.html}, \archivo{control_division.html}, \archivo{static/js/control.js}. & Romper AJAX, modales o captura de resultados. \\
Comunicaciones & \archivo{apps/invitations/services.py}, \archivo{forms.py}, \archivo{views.py}, comandos. & Enviar correos incorrectos o perder trazabilidad. \\
Django Admin & \archivo{apps/participants/admin.py}, \archivo{apps/tournament/admin.py}, \archivo{apps/invitations/admin.py}. & Ejecutar acciones operativas sin procedimiento. \\
\bottomrule
\end{longtable}

\section{Flujo recomendado para cambios}
\label{sec:mantenimiento-flujo-cambios}

Todo cambio debe seguir este flujo minimo:

\begin{enumerate}
  \item identificar modulo, flujo afectado y datos persistidos relacionados;
  \item revisar capitulos y anexos vinculados del manual;
  \item crear o actualizar pruebas si el cambio toca reglas de negocio;
  \item aplicar el cambio en una rama de trabajo;
  \item ejecutar \comando{manage.py check};
  \item ejecutar pruebas automaticas relacionadas en entorno no productivo;
  \item probar manualmente el flujo afectado;
  \item revisar que no se agreguen secretos, datos reales ni archivos generados;
  \item documentar la decision tecnica si cambia una regla del torneo.
\end{enumerate}

No se debe cambiar la estructura de modelos sin revisar migraciones y respaldos. No se debe cambiar \archivo{control.js} sin probar \variable{data-result-form}, modales, guardado masivo y refresco parcial.

\section{Uso de Git}
\label{sec:mantenimiento-git}

La auditoria no confirma una politica formal de versionado. Como practica minima, los cambios deben agruparse por proposito y evitar mezclar refactors con correcciones operativas. Antes de entregar una version se recomienda revisar:

\begin{itemize}
  \item \comando{git status};
  \item \comando{git diff --name-only};
  \item archivos generados o locales que no deben versionarse;
  \item migraciones nuevas cuando cambien modelos;
  \item documentacion asociada al cambio.
\end{itemize}

\section{Comandos y acciones operativas}
\label{sec:mantenimiento-comandos}

\begin{longtable}{p{3.8cm}p{4.8cm}p{4.2cm}}
\caption{Comandos y acciones con impacto operativo.}
\label{tab:mantenimiento-comandos}\\
\toprule
Accion & Uso & Precaucion \\
\midrule
\comando{bootstrap_tournament} & Crea edicion inicial y reglas base. & Ejecutar solo cuando se necesite inicializacion controlada. \\
\comando{reset_tournament_demo} & Limpia datos operativos y crea simulacion. & Destructivo; no usar en produccion. \\
\comando{send_invitations} & Procesa invitaciones desde XLSX. & Dry-run por defecto; envio real con \comando{--yes}. \\
\comando{send_attendance_requests} & Envia o simula solicitudes de asistencia. & Validar \variable{base-url}, SMTP y modo de envio. \\
Admin: \funcion{send_attendance_requests} & Accion sobre registros seleccionados. & Puede enviar correos; revisar seleccion y configuracion. \\
Admin: \funcion{sync_confirmed_to_teams} & Sincroniza registros confirmados a equipos. & Verificar confirmaciones antes de ejecutar. \\
\bottomrule
\end{longtable}

\section{Errores frecuentes}
\label{sec:mantenimiento-errores}

\begin{longtable}{p{4.2cm}p{4.5cm}p{4cm}}
\caption{Problemas frecuentes y respuesta recomendada.}
\label{tab:mantenimiento-errores}\\
\toprule
Problema & Causa probable & Respuesta \\
\midrule
Django no inicia en produccion & Falta \variable{DJANGO_SECRET_KEY}, \variable{DJANGO_ALLOWED_HOSTS} o variables PostgreSQL. & Revisar \archivo{.env} real y ejecutar \comando{manage.py check}. \\
Panel sin estilos o JS & Static files no recolectados o no servidos. & Ejecutar \comando{collectstatic} en entorno correcto y revisar servidor de estaticos. \\
No se guardan resultados AJAX & Token CSRF, atributos \variable{data-*} o respuesta JSON alterados. & Revisar \cref{chap:javascript-ajax} y probar formulario individual. \\
No se puede avanzar de fase & La fase no esta cerrada o faltan resultados. & Revisar \funcion{stage_is_closed} y resultados pendientes. \\
Correo real bloqueado & SMTP incompleto, placeholders o falta de confirmacion. & Validar estado de correo y usar prueba controlada. \\
PDF de invitacion falla & LibreOffice no disponible o timeout. & Configurar \variable{LIBREOFFICE_BINARY} y probar conversion. \\
\bottomrule
\end{longtable}

\section{Checklist antes de evento}
\label{sec:mantenimiento-checklist-evento}

Antes de operar el sistema en un evento real se recomienda:

\begin{itemize}
  \item validar \archivo{.env} sin placeholders;
  \item ejecutar \comando{manage.py check --deploy};
  \item confirmar base PostgreSQL y respaldo;
  \item probar SMTP en modo test;
  \item probar LibreOffice con una plantilla DOCX;
  \item crear edicion activa real;
  \item validar formularios de registro;
  \item confirmar usuarios staff autorizados;
  \item ejecutar simulacion completa de grupos, eliminatorias, repechaje, correccion manual y podio.
\end{itemize}

\section{Recomendaciones para futuras versiones}
\label{sec:mantenimiento-recomendaciones}

\begin{itemize}
  \item Separar \archivo{apps/tournament/services.py} por subdominios: inicializacion, resultados, fases manuales, recomendaciones, historial y sincronizacion.
  \item Documentar formalmente el schema de \variable{DivisionCompetition.configuration}.
  \item Crear runbooks para acciones de Admin y comandos de comunicaciones.
  \item Definir responsable de mantenimiento, politica de versionado y calendario de pruebas antes de evento.
  \item Agregar monitoreo de errores cuando exista entorno de produccion.
  \item Mantener el manual actualizado cuando cambien modelos, servicios o contratos AJAX.
\end{itemize}
